% Section 3.2, from lectures/L03/appendix/b_isr_contract.md.

\section{The ISR Contract, and Atomicity}
\label{c3:sec:contract}

\subsection{What the hardware saves}
\label{c3:sec:saves}
The program counter. That is the complete list.

Not \code{SREG}. Not any register. Not the pointer registers. Everything else your handler
touches, it destroys, and the interrupted code has no idea it was interrupted and will carry on
using whatever it had.

This is the difference between an interrupt handler and a subroutine, and it is not a small one. A
subroutine has a caller who agreed to a contract: \code{r18} to \code{r27}, \code{r30} and
\code{r31} are clobbered, and a caller who wanted one kept saved it first
(\secref{c2:sec:contract}). A handler has no caller. It interrupts code that agreed to nothing,
could not have agreed to anything, and does not know it happened.

\textbf{So a handler must save everything it uses, and restore all of it.} Not the call-saved
half: everything.

\subsection{\code{SREG} first, and \code{SREG} last}
\label{c3:sec:sreg}
The subtlest part is not a register at all.

The interrupted code may have been between a \code{cp} and its \code{brne}
(\secref{c1:sec:sreg}). The flags are live at that instant, in the strongest sense: they are the
entire result of the comparison, and they have one instruction left to live. A handler that
executes any arithmetic at all overwrites them, and the branch afterwards then tests the handler's
flags instead of the comparison's.

The symptom is a branch that goes the wrong way occasionally, and only under load, and never when
you are watching. It is among the least pleasant bugs this book can produce.

So a handler's prologue and epilogue have a fixed shape, and the order is not negotiable. On the
way in:
\begin{enumerate}
  \item \textbf{Push a register}, one you are willing to work in.
  \item \textbf{Read \code{SREG} into it} with \code{in}, before anything has had a chance to
    change the flags.
  \item \textbf{Push that too}, so the flags are on the stack.
\end{enumerate}
Then your handler does whatever it likes. On the way out, the mirror image:
\begin{enumerate}[resume]
  \item \textbf{Pop the saved flags} back into the register.
  \item \textbf{Write them to \code{SREG}} with \code{out}, before anything can disturb them
    again.
  \item \textbf{Pop the register} back to what the interrupted code had in it.
  \item \textbf{\code{reti}}.
\end{enumerate}

\textbf{Save the register before reading \code{SREG}}, because \code{push} writes no flags but you
need somewhere to put \code{SREG} and taking it costs you a register you have not saved yet.
\textbf{Restore \code{SREG} before the final \code{pop}}, for the same reason in reverse:
\code{pop} writes no flags either, so this order works, and an order that puts an arithmetic
instruction between the \code{out} and the \code{reti} does not.

\code{in} and \code{out} are the right instructions here: \code{SREG} is at data space
\code{0x5F}, which is I/O address \code{0x3F}, inside the low 64 that \code{in} and \code{out}
reach, and they cost one cycle where \code{lds} and \code{sts} cost two. Two cycles saved on every
interrupt, on the one path where it matters most, and two words of flash on top of that.

\subsection{Registers, and the cost of a handler}
\label{c3:sec:cost}
Beyond \code{SREG}, save every register your handler writes, including any that a subroutine you
call will clobber. That last part catches people: a handler that calls \code{btn_pressed} has to
save \code{r18}, \code{r19}, \code{r24} and \code{r25}, and the pointer registers \code{X} and
\code{Z} as well, even though it never names \code{r18}, \code{X} or \code{Z} itself, because the
subroutines it calls do.

Each \code{push} is 2 cycles and each \code{pop} is 2, so a saved register costs 4 cycles for the
round trip and one byte of stack. A handler saving eight registers plus \code{SREG} makes nine
pushes and nine pops: with the \code{in} and the \code{out} around them that is 19 cycles before it
does anything and 19 more on the way out, 38 in all.

That is the argument for short handlers, and it is worth being precise about what ``short'' buys
you, which is \secref{c3:sec:blocked}.

\subsection{The shared variable}
\label{c3:sec:shared}
A handler and the main loop that share a variable are two pieces of code with no synchronisation
between them, and on an 8-bit machine any variable wider than one byte is read and written in
pieces.

\bookfigure{atomicity}{Three moments in a row: the main loop reads the low byte of a counter
holding \code{0x00FF}, the interrupt lands and the handler increments the counter to
\code{0x0100}, and the main loop then reads the high byte, assembling
\code{0x01FF}.}{c3:fig:atomicity}

The main loop assembled \code{0x01FF}, which is 511. The counter was 255 and then 256. It was never
511, at any instant, and no amount of reading the two \code{lds} instructions will show you
anything wrong: each of them read the correct value of the byte it read, at the moment it read it.

\textbf{The window is exactly the gap between the two instructions}, which for two \code{lds} is 2
cycles. At 16\,MHz that is 125\,ns, and the chance that any one read is torn is that window times
the rate at which the handler actually fires: at 1\,kHz, about one read in eight thousand. Which
means it happens, and it happens rarely, and it will not happen while you are watching.

The fix is to make the read atomic by preventing the interrupt during it: \code{cli} before the two
loads and \code{sei} after them.

Two instructions and two cycles to close it. Note what \code{cli} and \code{sei} are doing here:
they are not protecting the \emph{counter}, they are protecting the \emph{pair of reads}. The
handler is free to change the counter at any other time.

\textbf{And note the bug this fix has.} That sequence unconditionally enables interrupts at the
end. If the code calling it had interrupts disabled for its own reasons, they are now on. The
correct version reads \code{SREG} into a spare register first, clears the flag, does the two loads,
and writes that register back to \code{SREG}; the same prologue and epilogue as
\secref{c3:sec:sreg}, for the same reason. It costs one cycle more, not zero: \code{cli} plus
\code{sei} is 2 cycles, and \code{in} plus \code{cli} plus \code{out} is 3. The \code{cli} does not
go away; only the \code{sei} does.

\textbf{And the failure that follows from the same bug, one level up.} Nest two of these. The outer
one does \code{cli}, and part-way through it calls something that does \code{cli}\,\dots\,\code{sei}.
The inner section ends by turning interrupts \emph{on}, in the middle of the outer section, which
now runs unprotected to its end and has no way to find out. The \code{sei} version is not merely
impolite to its caller; it is impossible to compose with itself. Save-and-restore composes, because
each level hands back the flag it was given, and that is the whole reason a kernel's critical
sections cannot be \code{cli}/\code{sei} either.

\textbf{A single byte needs none of this.} \code{lds r24, flag} is one instruction, and an
interrupt cannot land in the middle of one. That is why an 8-bit flag shared between a handler and
a main loop is safe and a 16-bit counter is not, and it is worth knowing which of your variables is
which.

\subsection{How long are interrupts off}
\label{c3:sec:blocked}
The hardware clears \code{I} before it pushes the program counter and \code{reti} sets it again.
So the window in which no other interrupt can be served is:
\[
  \text{blocked} = 4 \;+\; \text{vector jump} \;+\; \text{your handler, \code{reti} included}
\]
Four for the push, two more for an \code{rjmp} in the table, and then everything you wrote. Six
cycles are the hardware's and the rest is yours.

On the reference implementation, the button handler measures \textbf{104 cycles when the button
is not pressed and 163 when it is}, so it blocks interrupts for \textbf{110} or \textbf{169} cycles,
which is 6.875\,µs or 10.5625\,µs at 16\,MHz. Whether that is a lot depends entirely on what else is
waiting; the cross-check in Exercise~\ref{c3:ex:crosscheck} is where you work it out for your own
handler and check the answer against a measurement.

\textbf{What being blocked costs is lateness, not loss.} A pin change interrupt sets a flag, and
the flag stays set until it is served, so a second interrupt arriving during your handler is
delayed rather than dropped. What \emph{is} lost is a second change on the same port during that
window: the flag is one bit, and it is already set.

\subsection{What a handler must not do}
\label{c3:sec:mustnot}
\textbf{Do not call anything long.} Not because it is bad style, but because
\secref{c3:sec:blocked} is a number and calling a delay routine makes it enormous.

\textbf{Do not busy-wait for anything an interrupt would have to deliver.} Interrupts are off. The
thing you are waiting for cannot arrive, and the machine will sit there until the watchdog
(Chapter~6) rescues it, or until somebody unplugs it.

\textbf{Do not \code{sei} inside a handler}, unless you have thought about it very hard. It lets
the same interrupt interrupt itself, and the stack grows by one frame each time. Two frames is
fine. Two thousand frames is your variables overwritten from the top down, with no fault and no
message (\secref{c2:sec:stack}).

\textbf{Do not assume it runs once per event.} A mechanical button bounces
(\secref{c2:sec:button}), and a pin change interrupt fires on every one of those transitions. A
handler that toggles an LED on each press will run four to ten times per press, and the LED will
appear to respond about half the time. That is not a fault in your handler and no amount of
rereading it will help.

\textbf{And do not forget it can happen between any two instructions of anything.} Including
between the two \code{lds} of \secref{c3:sec:shared}, and including inside a subroutine the main
loop was halfway through.
